一个自旋 \(1/2\) 粒子初态为 \(|+z\rangle\)。先测量沿单位方向 \(\boldsymbol n=(\sin\theta,0,\cos\theta)\)（\(0<\theta<\pi\)）的自旋投影 \(\boldsymbol n\cdot\boldsymbol\sigma\)，其本征值为 \(\pm1\)。此后粒子在均匀磁场中按哈密顿量 \(H=-\hbar\omega\sigma_z/2\) 演化时间 \(t\ge0\)，\(\omega>0\)，再测量 \(\sigma_x\)。
(1) 求首次测到 \(+1\) 的概率及相应归一化后测态；(2) 在选择首次 \(+1\) 结果的子样本中，求末次测到 \(\sigma_x=+1\) 的条件概率；(3) 若首次测量仍实施但不记录结果，利用密度矩阵求末次测量概率，并与未经首次测量的情况比较；(4) 对 \(\theta=\pi/2\) 给出三个概率随时间的差异，说明为何丢弃测量记录不能等价于未作测量。
